\section{The Spine* Algorithm}
\label{sec:algorithm}

\subsection{Preliminaries}

Let $G = (V, E, w)$ be a planar graph representing the census units (tracts
or blocks) of a state, with vertex weights $w_v$ equal to the 2020 census
population of unit $v$.  Let $C$, $S$, and $H$ be the state's congressional,
senate, and house seat counts under the 2020 apportionment.

A \emph{bisection spine} for a chamber with $k$ seats is a sequence of
positive integers $[s_1, s_2, \ldots, s_\ell]$ with $\prod_i s_i = k$,
interpreted as a recursion schedule: first split $G$ into $s_1$ sub-graphs,
then each sub-graph into $s_2$ sub-sub-graphs, and so on.  The canonical
choice is the prime factorization of $k$ in non-decreasing order.

\begin{definition}[Compatible Spines]
\label{def:compatible}
A \emph{compatible triple of spines} for $(C, S, H)$ is a triple
$(P_C, P_S, P_H)$ of bisection spines such that $P_C$, $P_S$, and $P_H$
share a common prefix $\tau$, called the \emph{trunk}.  The trunk product
$\prod \tau$ equals $g = \gcd(C, S, H)$.
\end{definition}

The trunk $\tau$ is the prime factorization of $g$ in non-decreasing order.
The chamber-specific tails are the prime factorizations of $C/g$, $S/g$,
and $H/g$ respectively.  A tail of length zero (quotient $= 1$) means the
chamber ends exactly at the trunk level.

\subsection{Algorithm: \textsc{CompatibleSpines}}

\begin{algorithm}
\caption{\textsc{CompatibleSpines}$(C, S, H)$}
\label{alg:compatible_spines}
\begin{algorithmic}[1]
\Require $C, S, H \geq 1$
\Ensure $(P_C, P_S, P_H)$ --- compatible triple of bisection spines
\State $g \leftarrow \gcd(C, S, H)$
\State $\tau \leftarrow \textsc{PrimeFactors}(g)$ \Comment{sorted non-decreasing}
\State $\text{tail}_C \leftarrow \textsc{PrimeFactors}(C / g)$
\State $\text{tail}_S \leftarrow \textsc{PrimeFactors}(S / g)$
\State $\text{tail}_H \leftarrow \textsc{PrimeFactors}(H / g)$
\State $P_C \leftarrow \tau \Vert \text{tail}_C$
\State $P_S \leftarrow \tau \Vert \text{tail}_S$
\State $P_H \leftarrow \tau \Vert \text{tail}_H$
\State \Return $(P_C, P_S, P_H)$
\end{algorithmic}
\end{algorithm}

\begin{proposition}[Correctness]
\label{prop:correct}
For any $C, S, H \geq 1$, $\textsc{CompatibleSpines}(C, S, H)$ returns
a triple $(P_C, P_S, P_H)$ satisfying:
\begin{enumerate}
  \item $\prod P_C = C$, $\prod P_S = S$, $\prod P_H = H$.
  \item $P_C$, $P_S$, $P_H$ all begin with $\tau = \textsc{PrimeFactors}(\gcd(C,S,H))$.
  \item The trunk $\tau$ is the \emph{longest} common prefix that can be
        guaranteed for any $(C, S, H)$ with the given GCD.
\end{enumerate}
\end{proposition}

\begin{proof}
(1) The product of a spine equals the product of its trunk and tail.
$\prod \tau = g$ and $\prod \text{tail}_C = C/g$, so $\prod P_C = g \cdot (C/g) = C$.
Analogously for $S$ and $H$.

(2) By construction, all three spines begin with $\tau$.

(3) Suppose a longer common prefix $\tau'$ with $\prod \tau' = g' > g$ existed.
Then $g'$ would divide all of $C$, $S$, $H$, contradicting $g = \gcd(C,S,H)$.
\end{proof}

\subsection{Compatibility Score}

\begin{definition}[Spine Compatibility Score]
\label{def:score}
The \emph{spine compatibility score} of $(C, S, H)$ is:
\begin{equation}
  \sigma(C, S, H) \;=\; \left(1 - \frac{\gcd(C,S,H)}{\min(C,S,H)}\right) \times 100.
\end{equation}
\end{definition}

\begin{lemma}[Score Properties]
\label{lem:score}
\begin{enumerate}
  \item $\sigma(C, S, H) \in [0, 100]$.
  \item $\sigma(C, S, H) = 0$ if and only if $\gcd(C,S,H) = \min(C,S,H)$, i.e.,
        the smallest chamber's count divides both larger counts.
  \item $\sigma(C, S, H) = 100$ only in the degenerate limit; the maximum
        attained score for any US state is $\approx 96.8$ (Texas: $g=1$,
        $\min = 31$).
\end{enumerate}
\end{lemma}

\begin{proof}
(1) Since $\gcd(C,S,H) \leq \min(C,S,H)$, the fraction is in $[0,1]$.
(2) Immediate from the definition.
(3) The score equals $100(1 - 1/m)$ when $g=1$.  For any US state,
$\min(C,S,H) \leq C \leq 52$, so the score is at most $100(1 - 1/52) \approx 98.1$,
but no state achieves this because $C=52$ (California) has $g=4 > 1$.
\end{proof}

\begin{theorem}[Bimodality Gap]
\label{thm:bimodal}
For positive integers $C$, $S$, $H$ with $g = \gcd(C,S,H)$ and
$m = \min(C,S,H)$, the compatibility score $\sigma = (1 - g/m)\times 100$
satisfies
\[
  \sigma \;\in\; \{0\} \;\cup\; [50, 100).
\]
There are no states with $0 < \sigma < 50$.
\end{theorem}

\begin{proof}
Suppose for contradiction that $0 < \sigma < 50$.  Then
\[
  0 \;<\; 1 - \frac{g}{m} \;<\; \frac{1}{2},
  \quad\text{i.e.,}\quad
  \frac{m}{2} \;<\; g \;<\; m.
\]
Since $g = \gcd(C,S,H)$ divides every element, in particular $g \mid m$.
We therefore need a divisor of $m$ that lies strictly between $m/2$ and $m$.
Such a divisor would be $m/p$ for some integer $p$ with $1 < p < 2$ ---
but no such integer exists, since no prime is less than~$2$.

More precisely: if $g \mid m$ and $g > m/2$, then $m/g < 2$, so $m/g = 1$
(it must be a positive integer), which gives $g = m$ and $\sigma = 0$ ---
a contradiction with $\sigma > 0$.

\medskip
\noindent\textit{Boundary cases.}
\begin{itemize}
  \item If $m = 1$, then $g = \gcd(C,S,H) = 1$ (the only positive divisor of
        $1$ is $1$ itself), so $\sigma = (1 - 1/1)\times 100 = 0$.  Consistent.
  \item The infimum of the non-zero range is achieved when $g = 1$ and $m = 2$:
        $\sigma = (1 - 1/2)\times 100 = 50$.  Hence the non-zero values begin
        at exactly~$50$, confirming $\sigma \in \{0\} \cup [50, 100)$.
\end{itemize}
\end{proof}

\begin{remark}
The Bimodality Gap theorem is a pure number-theoretic fact: it holds for any
positive integers $(C, S, H)$, not merely for current US apportionments.
The empirical observation that no US state falls in the range $0 < \sigma < 50$
is therefore not an accident of the current apportionment cycle but a structural
invariant of the GCD-based score.
\end{remark}

\subsection{Nesting Modes}

The compatibility score determines which nesting mode is achievable.

\paragraph{Mode 1 --- Strict nesting ($\sigma = 0$).}
When $\gcd(C,S,H) = \min(C,S,H)$, the smallest chamber is exactly the trunk.
Its districts are the trunk regions.  The larger chambers are refinements of
the trunk.  No boundary violations occur.

\paragraph{Mode 2 --- Partial nesting ($0 < \sigma < 50$).}
The trunk covers more than half the structural degrees of freedom of the
smallest chamber.  The spine provides a strong geographic anchor, but the
tail splits produce boundary zones where chambers partially overlap.

\paragraph{Mode 3 --- Best-effort nesting ($\sigma \geq 50$).}
The trunk is a small fraction of the smallest chamber.  The chambers can
share only a coarse geographic skeleton.  A mismatch tolerance $\tau_\text{pop}$
(population fraction allowed in boundary zones) is required.

\subsection{The Full NestSection Algorithm}

\begin{algorithm}
\caption{\textsc{NestSection}$(G, C, S, H, \tau_\text{pop})$}
\label{alg:nestsection}
\begin{algorithmic}[1]
\Require Graph $G$, seat counts $C, S, H$, tolerance $\tau_\text{pop} \geq 0$
\Ensure Tri-chamber district plan $(D_C, D_S, D_H)$
\State $(P_C, P_S, P_H) \leftarrow \textsc{CompatibleSpines}(C, S, H)$
\State $g \leftarrow \gcd(C, S, H)$
\State $\mathcal{R} \leftarrow \textsc{GeoSection}(G, g)$
  \Comment{$g$ trunk regions via GeoSection~\citep{B8}}
\ForAll{trunk region $R_i \in \mathcal{R}$, $i = 1, \ldots, g$}
  \State $c_i \leftarrow \textsc{ApportionRegions}(\{w(R_j)\}_j, C)[i]$
    \Comment{Integer seats per trunk~\citep{B11}}
  \State $s_i \leftarrow S / g$
  \State $h_i \leftarrow H / g$
  \State $D_C[i] \leftarrow \textsc{GeoSection}(R_i, c_i)$
  \State $D_S[i] \leftarrow \textsc{GeoSection}(R_i, s_i)$
  \State $D_H[i] \leftarrow \textsc{NestRefine}(D_S[i], h_i / s_i, \tau_\text{pop})$
\EndFor
\State \Return $(\bigcup_i D_C[i],\; \bigcup_i D_S[i],\; \bigcup_i D_H[i])$
\end{algorithmic}
\end{algorithm}

\paragraph{NestRefine.}  \textsc{NestRefine} splits each senate district into
$h_i / s_i$ house sub-districts.  When $h_i / s_i$ is not an integer, it
assigns floor or ceiling counts by population weight and marks boundary zones
where a house district straddles a senate boundary.  The population fraction
in boundary zones must not exceed $\tau_\text{pop}$.

\subsection{Trunk Size and Geographic Scale}

The trunk size $g$ determines the geographic scale of the shared structure.
For Oregon ($g = 6$), the shared trunk has 6 regions averaging roughly 700,000
residents each --- approximately the size of a single congressional district.
For Alabama ($g = 7$), the trunk has 7 regions matching the 7 congressional
districts exactly, and senate districts are unions of exactly 5 trunk regions
each.

For states with $g = 1$ (e.g., Texas with $C=38$, $S=31$, $H=150$), the
``trunk'' is the entire state, and the three chambers share no sub-state
geographic anchor beyond the state boundary.  In these cases NestSection
reduces to three independent GeoSection calls with the shared property that
they start from the same census graph $G$ --- a weak form of consistency.

\subsection{Complexity}

Algorithm~\ref{alg:compatible_spines} runs in $O(\log(\max(C,S,H)))$ time
(GCD computation plus factorization).  The dominant cost of
Algorithm~\ref{alg:nestsection} is the $g+3$ calls to \textsc{GeoSection},
each of which runs in time proportional to graph-partitioning cost.  The
total additional cost over three independent GeoSection runs is the $g$
trunk-level calls, which is at most $\min(C,S,H)$ additional partitioning
operations.
